\documentclass[11pt,a4paper]{article}
\usepackage[T1]{fontenc}
\usepackage[utf8]{inputenc}
\usepackage[brazil]{babel}
\usepackage{lmodern}
\usepackage[a4paper,margin=2.2cm]{geometry}
\usepackage{graphicx}
\usepackage{booktabs}
\usepackage{xcolor}
\usepackage{tikz}
\usetikzlibrary{positioning,arrows.meta,calc,fit,backgrounds}
\usepackage{caption}
\usepackage{subcaption}
\usepackage{enumitem}
\usepackage[hidelinks]{hyperref}
\usepackage{float}
\usepackage{amsmath}
\usepackage{newunicodechar}
\usepackage{microtype}
\usepackage{titlesec}
\usepackage{fancyhdr}
\newunicodechar{√}{\ensuremath{\surd}}
\newunicodechar{м}{m}\newunicodechar{о}{o}\newunicodechar{а}{a}\newunicodechar{т}{t}

\definecolor{raia}{HTML}{0B3D91}
\definecolor{ciano}{HTML}{1AA7B5}
\captionsetup{font=small,labelfont={bf,color=raia},format=hang}
\titleformat{\section}{\large\bfseries\color{raia}}{\thesection}{0.6em}{}[{\color{raia!35}\titlerule[0.6pt]}]
\titleformat{\paragraph}[runin]{\normalsize\bfseries\color{raia}}{}{0pt}{}
\pagestyle{fancy}\fancyhf{}
\fancyhead[L]{\small\color{raia!70}RAIA: natação universitária InterUSP e TUSCA}
\fancyhead[R]{\small\color{raia!70}SCC5836, 1ª entrega}
\fancyfoot[C]{\small\thepage}
\renewcommand{\headrulewidth}{0.4pt}
\setlength{\headheight}{14pt}
\newcommand{\moldura}[2][\linewidth]{{\setlength{\fboxsep}{2pt}\setlength{\fboxrule}{0.7pt}%
  \fcolorbox{raia!40}{white}{\includegraphics[width=\dimexpr#1-5.4pt\relax]{#2}}}}
\setlist{itemsep=1pt,topsep=3pt}
\graphicspath{{figs/}}

\title{\textbf{RAIA}: exploração visual da natação universitária InterUSP e TUSCA (2023--2026)\\[4pt]
\large Projeto de Visualização Computacional --- 1ª entrega}
\author{Jáder Louis de Souza Gonçalves \quad Nº USP 17753116\\
\small SCC5836 --- Visualização Computacional --- ICMC-USP --- 2026/2}
\date{15 de novembro de 2026}


\definecolor{destaque}{HTML}{FFD23F}
\newcommand{\anotada}[3][\linewidth]{{\setlength{\fboxsep}{2pt}\setlength{\fboxrule}{0.7pt}%
  \fcolorbox{raia!40}{white}{\begin{tikzpicture}
    \node[inner sep=0, anchor=south west] (img) {\includegraphics[width=\dimexpr#1-5.4pt\relax]{#2}};
    \begin{scope}[x={(img.south east)}, y={(img.north west)}]#3\end{scope}
  \end{tikzpicture}}}}
\newcommand{\destaque}[5]{%
  \draw[destaque, line width=1.1pt, rounded corners=1.5pt] ({#1},{1-#2}) rectangle ({#3},{1-#4});
  \node[circle, fill=destaque, draw=raia, line width=0.4pt, text=raia, font=\sffamily\bfseries\scriptsize,
        inner sep=0pt, minimum size=3.6mm] at ({#1},{1-#2}) {#5};}
\DeclareRobustCommand{\marca}[1]{\tikz[baseline=-0.6ex]\node[circle, fill=destaque, draw=raia, line width=0.3pt, text=raia,
  font=\sffamily\bfseries\tiny, inner sep=0pt, minimum size=2.9mm]{#1};}

\begin{document}
\maketitle
\thispagestyle{empty}

\begin{abstract}
A RAIA reúne os resultados de natação das duas maiores competições universitárias paulistas, o InterUSP e o TUSCA,
entre 2023 e 2026, num único conjunto de dados com 142 provas, 1.816 resultados e 504 atletas. O conjunto foi
construído a partir das súmulas oficiais e é explorado num sistema web com uma piscina 3D e visualizações
interativas em d3. Esta primeira entrega descreve os dados, o pré-processamento e quatro questões de investigação,
respondidas com sete técnicas de visualização. O sistema está em \url{https://raia.jaderlouis.com.br}; o código e os dados, em \url{https://github.com/Prism411/raia-scc5836}; e o vídeo de apresentação, em \url{https://www.youtube.com/watch?v=95Qi5o8GiNw}.
\end{abstract}

\section{Introdução}\label{sec:intro}

\paragraph{Contextualização.} A natação competitiva segue as regras da World Aquatics~\cite{wa-regras}, adotadas
no Brasil pela Confederação Brasileira de Desportos Aquáticos (CBDA)~\cite{cbda} e, em São Paulo, pela Federação
Aquática Paulista (FAP)~\cite{fap}. No esporte universitário, além dos Jogos Universitários Brasileiros da
CBDU~\cite{cbdu}, as competições entre atléticas reúnem centenas de nadadores por ano. Duas das maiores do estado
são o InterUSP, disputado entre as atléticas da USP, e o TUSCA, o torneio universitário de São Carlos. Ambos usam
as provas oficiais (distâncias e estilos) e publicam súmulas com colocação, tempo e pontos de cada nadador.

\paragraph{Problema.} As súmulas saem como PDFs isolados, uma prova por arquivo e com um layout por edição. Com
elas não dá para responder perguntas simples: o nível está subindo? As duas competições são parecidas? Que equipe
domina cada prova? Onde um nadador perde tempo? Também não existe um conjunto de dados consolidado dessas
competições.

\paragraph{Motivação.} Como nadador do TUSCA, convivo com essas perguntas e com a dificuldade de respondê-las.
Atletas, técnicos e organizadores tomam decisões (em que prova nadar, o que treinar, como montar uma equipe) sem
uma visão de conjunto dos resultados.

\paragraph{Proposta.} A RAIA transforma as súmulas de 2023 a 2026 num conjunto único e o explora de duas formas:
uma \emph{análise visual}, com visões interativas sobre nível técnico, equipes e fases da prova, e uma
\emph{experiência}, em que o visitante monta um nadador e nada uma final real numa piscina 3D. A análise responde
às questões da Seção~\ref{sec:dados}; a experiência aproxima o público das competições.

\paragraph{Organização do texto.} A Seção~\ref{sec:dados} descreve o conjunto de dados e as questões de
investigação; a Seção~\ref{sec:prep}, o pré-processamento; a Seção~\ref{sec:analise}, as visualizações e o que
elas mostram; a Seção~\ref{sec:disc}, as discussões preliminares e as limitações; e a Seção~\ref{sec:prox}, os
próximos passos.

\section{Conjunto de dados e questões}\label{sec:dados}

\paragraph{Conjunto de dados.} Os resultados vêm das páginas públicas do EBAN Swim, que publica as súmulas oficiais
do InterUSP~\cite{eban-interusp} e do TUSCA~\cite{eban-tusca}. A Tabela~\ref{tab:dados} resume o conjunto.
Cada resultado tem colocação, atleta, equipe, tempo, pontos, categoria (oficial ou extra) e situação
(válido, desclassificado ou não terminou); cada prova tem competição, ano, sexo, distância, estilo, se é
revezamento ou demonstração, recorde vigente e vencedor.

\begin{table}[H]
\centering\small
\begin{tabular}{lll}
\toprule
Entidade & Quantidade & Variáveis \\
\midrule
Provas (finais) & 142 & competição, ano (2023--2026), sexo, distância, estilo, revezamento, recorde \\
Resultados & 1.816 & colocação, tempo (s), pontos, categoria, situação \\
Atletas únicos & 504 & nome, sexo, equipe, competições, anos (45 nadaram nas duas) \\
Equipes & 19 & 9 no InterUSP e 12 no TUSCA (POLI e CAASO nas duas) \\
\bottomrule
\end{tabular}
\caption{O conjunto de dados. Variáveis categóricas (competição, equipe, estilo, sexo), ordinais (ano,
colocação) e quantitativas (tempo, pontos, Índice Técnico).}
\label{tab:dados}
\end{table}

Como tempos de provas diferentes não são comparáveis, uso uma variável derivada, o \textbf{Índice Técnico (IT)}:
\[
\mathrm{IT} = 1000 \left(\frac{t_{\text{ref}}}{t}\right)^3,
\]
em que $t$ é o tempo do resultado e $t_{\text{ref}}$ é a melhor marca da base para a mesma prova e sexo. A forma
cúbica segue a pontuação da World Aquatics~\cite{wa-points}, com a melhor marca da base no lugar do recorde mundial:
o melhor resultado vale 1000 e cada 1\% a mais de tempo custa cerca de 3\% do índice.

\paragraph{Questões.} Uma análise preliminar das súmulas levou a quatro questões de investigação:
\begin{enumerate}[label=\textbf{Q\arabic*},leftmargin=2.2em]
  \item O nível técnico está subindo ao longo das edições?
  \item O InterUSP e o TUSCA têm o mesmo nível técnico?
  \item Que equipes dominam cada prova, e isso muda com o tempo?
  \item Onde um nadador perde tempo em relação à elite: na saída, nas viradas ou no nado?
\end{enumerate}

\section{Pré-processamento}\label{sec:prep}

As súmulas foram baixadas das páginas de resultados; os resultados do InterUSP 2023--2025 foram compilados numa
planilha, e o InterUSP 2026 e o TUSCA 2023--2026 vêm em PDF (108 arquivos). O script \texttt{dados/construir.py}
(Python) transforma essas fontes no arquivo \texttt{natacao.json} que a interface lê, em sete etapas
(Figura~\ref{fig:pipeline}).

\begin{figure}[H]
\centering
\resizebox{\linewidth}{!}{%
\begin{tikzpicture}[
  font=\sffamily, >={Stealth[length=2.2mm]},
  ic/.style={circle, draw=raia, line width=0.8pt, fill=raia!7, minimum size=12mm, inner sep=0pt},
  num/.style={circle, fill=ciano, text=white, font=\sffamily\bfseries\tiny, minimum size=3.6mm, inner sep=0pt},
  tit/.style={font=\sffamily\bfseries\scriptsize, text=raia, align=center},
  sub/.style={font=\sffamily\tiny, text=black!60, align=center, text width=17.5mm},
  seta/.style={->, draw=raia!70, line width=0.9pt},
  painel/.style={rounded corners=4pt, draw=raia!25, fill=raia!3, line width=0.6pt}
]
\tikzset{
  pics/doc/.style={code={
    \draw[raia, line width=0.6pt, fill=white] (-0.17,-0.22) -- (-0.17,0.22) -- (0.08,0.22) -- (0.17,0.13) -- (0.17,-0.22) -- cycle;
    \draw[raia, line width=0.5pt] (0.08,0.22) -- (0.08,0.13) -- (0.17,0.13);
    \foreach \y in {0.05,-0.03,-0.11} \draw[raia!60, line width=0.5pt] (-0.11,\y) -- (0.11,\y);
    \node[font=\sffamily\bfseries\fontsize{3.5}{4}\selectfont, text=white, fill=#1, inner sep=0.6pt, rounded corners=0.5pt] at (0,-0.16) {\tikzpictext};
  }},
  pics/texto/.style={code={\node[font=\sffamily\bfseries\large, text=raia] at (0,0) {Á\kern-0.5pt a};}},
  pics/cadeado/.style={code={
    \draw[raia, line width=1pt] (-0.1,0.02) -- (-0.1,0.12) arc (180:0:0.1) -- (0.1,0.02);
    \draw[raia, fill=raia, rounded corners=1pt] (-0.16,-0.2) rectangle (0.16,0.04);
    \fill[white] (0,-0.07) circle (0.035); \fill[white] (-0.012,-0.15) rectangle (0.012,-0.07);
  }},
  pics/ajuste/.style={code={
    \foreach \y/\x in {0.12/-0.05, 0/0.08, -0.12/-0.1} {
      \draw[raia!55, line width=0.7pt] (-0.19,\y) -- (0.19,\y);
      \fill[raia] (\x,\y) circle (0.045);
    }
  }},
  pics/rede/.style={code={
    \coordinate (p1) at (-0.15,0.1); \coordinate (p2) at (0.15,0.12); \coordinate (p3) at (0,-0.14);
    \draw[raia!60, line width=0.7pt] (p1) -- (p2) -- (p3) -- cycle;
    \foreach \p in {p1,p2,p3} \fill[raia] (\p) circle (0.055);
  }},
  pics/barras/.style={code={
    \draw[raia!60, line width=0.5pt] (-0.19,-0.17) -- (0.19,-0.17);
    \fill[raia!60] (-0.15,-0.17) rectangle (-0.07,-0.02);
    \fill[raia] (-0.04,-0.17) rectangle (0.04,0.17);
    \fill[ciano] (0.07,-0.17) rectangle (0.15,0.07);
  }},
  pics/confere/.style={code={
    \draw[ciano, line width=1.6pt, line cap=round, line join=round] (-0.13,0) -- (-0.03,-0.11) -- (0.15,0.12);
  }},
  pics/tela/.style={code={
    \draw[raia, line width=0.7pt, fill=white, rounded corners=1pt] (-0.24,-0.12) rectangle (0.24,0.2);
    \fill[ciano!35] (-0.2,-0.08) rectangle (0.2,0.16);
    \foreach \y in {-0.02,0.05,0.12} \draw[white, line width=0.5pt, dash pattern=on 1pt off 1pt] (-0.2,\y) -- (0.2,\y);
    \draw[raia, line width=0.7pt] (0,-0.12) -- (0,-0.19); \draw[raia, line width=0.7pt] (-0.09,-0.19) -- (0.09,-0.19);
  }}
}
\node[painel, minimum width=20mm, minimum height=34mm] (fontes) at (0,0) {};
\node[tit, anchor=north] at (fontes.north) [yshift=-1.2mm] {Fontes};
\pic[pic text=PDF, scale=1.6] at (-0.4,0.42) {doc=red!70!black};
\pic[pic text=CSV, scale=1.6] at (0.4,0.42) {doc=green!45!black};
\node[sub, anchor=north] at (0,-0.1) {108 súmulas PDF\\e a planilha do\\InterUSP 2023--25\\(EBAN Swim)};
\foreach \i/\icone/\t/\s in {
  1/doc/Extração/{pdftotext e\\expressões regulares},
  2/texto/Codificação/{acentos quebrados\\e letras cirílicas},
  3/cadeado/Privacidade/{CPFs removidos\\na extração},
  4/ajuste/Normalização/{equipes, tempos\\e categorias},
  5/rede/Entidades/{504 atletas\\ligados entre anos},
  6/barras/Variáveis/{vencedor, recorde\\e nº de nadadores},
  7/confere/Validação/{65 de 71 totais\\conferidos}} {
  \node[ic] (e\i) at ({2.5+1.9*(\i-1)},0.32) {};
  \ifnum\i=1 \pic[pic text=PDF] at (e\i) {doc=raia};
  \else \pic at (e\i) {\icone}; \fi
  \node[num] at ($(e\i)+(-0.43,0.43)$) {\i};
  \node[tit, below=1.2mm of e\i] (t\i) {\t};
  \node[sub, below=0mm of t\i] {\s};
}
\foreach \a/\b in {1/2,2/3,3/4,4/5,5/6,6/7} \draw[seta] (e\a) -- (e\b);
\draw[seta] (fontes.east |- e1) -- (e1);
\begin{scope}[on background layer]
  \node[painel, fill=raia!2, fit=(e1)(e7)(t1)(t7), inner xsep=2mm, inner ysep=7mm, yshift=-1.5mm] (faixa) {};
\end{scope}
\node[font=\sffamily\bfseries\tiny, text=raia!70, anchor=south west] at (faixa.north west) {dados/construir.py};
\node[painel, minimum width=22mm, minimum height=34mm, right=5mm of faixa.east, anchor=west] (saida) {};
\node[tit, anchor=north] at (saida.north) [yshift=-1.2mm] {Saída};
\pic[pic text=JSON, scale=1.6] at ($(saida.center)+(-0.5,0.42)$) {doc=raia};
\pic[scale=1.4] at ($(saida.center)+(0.48,0.42)$) {tela};
\draw[seta, line width=0.6pt] ($(saida.center)+(-0.18,0.42)$) -- ($(saida.center)+(0.08,0.42)$);
\node[sub, anchor=north, text width=21mm] at ($(saida.center)+(0,-0.1)$) {natacao.json\\(142 provas, 1.816\\resultados) lido\\pela RAIA};
\draw[seta] (faixa.east |- e7) -- (saida.west |- e7);
\end{tikzpicture}}
\caption{Do PDF à visualização: as sete etapas de \texttt{dados/construir.py}, que transformam as súmulas no
\texttt{natacao.json} lido pela interface.}
\label{fig:pipeline}
\end{figure}

\begin{itemize}[leftmargin=1.2em]
  \item \textbf{Extração.} Cada PDF é convertido em texto com \texttt{pdftotext}; expressões regulares reconhecem o
  título da prova e cada linha de resultado. A equipe é identificada no fim de cada linha a partir de uma lista
  de grafias conhecidas, o que resolve linhas sem espaçamento regular.
  \item \textbf{Codificação.} Alguns PDFs trazem acentos quebrados (``JO√O''), caracteres invisíveis e letras
  cirílicas visualmente idênticas às latinas no meio de nomes, que são corrigidos.
  \item \textbf{Privacidade.} As súmulas do TUSCA 2026 trazem o CPF dos atletas. Os CPFs são removidos logo após a
  extração, o script falha se sobrar algum na saída, e os PDFs brutos ficam fora do repositório.
  \item \textbf{Normalização.} Grafias da mesma equipe são unificadas (MEDICINA PINHEIROS, MED.PINHEIROS e
  PINHEIROS viram MED PINHEIROS), tempos viram segundos (``1.02.45'' $\to$ 62{,}45) e a categoria (oficial ou
  extra) é lida da súmula ou inferida dos pontos.
  \item \textbf{Entidades.} O mesmo atleta é ligado entre edições e competições: mesmo sexo e primeiro nome e,
  na mesma equipe, último nome a no máximo uma edição de distância (Levenshtein) ou um nome contido no outro.
  \item \textbf{Variáveis.} Vencedor, número de nadadores, recorde vigente e quebra de recorde por prova; o IT e o
  modelo de fases da prova (Seção~\ref{sec:analise}) são calculados na interface.
  \item \textbf{Validação.} A soma dos pontos por equipe confere com a classificação geral das súmulas em 65 de 71
  linhas; as diferenças vêm de critérios de desempate e de bônus de recorde. Duas súmulas tinham colocações
  digitadas fora de ordem, refeitas pela ordem dos tempos.
\end{itemize}

\section{Análise visual}\label{sec:analise}

O sistema é uma página web única (React, d3~\cite{d3} e three.js~\cite{threejs}). Na análise, os filtros por
competição, ano, sexo, prova e equipe, a cor por equipe, competição ou estilo, as dicas ao passar o mouse e a
ficha de cada atleta ao clicar seguem o mantra ``visão geral primeiro, zoom e filtro, detalhes sob
demanda''~\cite{shneiderman}. As cores seguem a paleta de Okabe e Ito~\cite{okabe-ito}, segura para daltonismo, e um
simulador de daltonismo no topo da página aplica protanopia, deuteranopia, tritanopia ou acromatopsia à página
inteira. A escolha de cada técnica parte do tipo de dado e da tarefa~\cite{munzner}; a Tabela~\ref{tab:visoes} resume as
técnicas e a questão que cada uma responde.

\begin{table}[H]
\centering\small
\begin{tabular}{p{3.3cm}p{4.4cm}p{1.1cm}p{5.6cm}}
\toprule
Visão & Técnica & Questão & Por que essa técnica \\
\midrule
Nível & \emph{ridgeline} de densidades & Q1 & compara a distribuição inteira do IT entre 8 edições, sem perder a forma \\
Enxame & \emph{beeswarm} (um ponto por resultado) & Q2 & mostra cada resultado e os \emph{outliers}, sem esconder a dispersão numa média \\
Ranking & \emph{bump chart} & Q3 & posição ordinal das equipes ao longo dos anos \\
Equipe $\times$ prova & matriz (mapa de calor) reordenável & Q3 & duas variáveis categóricas e uma quantitativa; a reordenação revela blocos~\cite{behrisch} \\
Raio-X & barras empilhadas no tempo (fita), linhas e faixas & Q4 & decompõe o tempo da prova em fases \\
Arena 3D & animação da final em 3D & Q4 & mostra a prova como ela acontece, com placar e telemetria \\
Ficha do atleta & linha (trajetória) e radar & --- & detalhe sob demanda de um atleta \\
\bottomrule
\end{tabular}
\caption{Técnicas de visualização da 1ª entrega.}
\label{tab:visoes}
\end{table}

\paragraph{Q1 e Q2: nível técnico.} O \emph{ridgeline} (Figura~\ref{fig:nivel}) empilha a densidade do IT de cada
edição, com a mediana tracejada. No InterUSP a mediana caiu de 556 (2023) para 512 (2025) e voltou a 524 (2026);
no TUSCA caiu de 633 para 506 entre 2023 e 2024 e ficou estável depois (516 e 524). As densidades mostram que
as duas quedas são diferentes. No InterUSP as finais cresceram (de 194 para 216 resultados) e o meio da
distribuição desceu, mas a cauda direita subiu: o percentil 90 do IT foi de 802 para 853. No TUSCA a
distribuição inteira desceu em 2024, inclusive o topo (percentil 90 de 876 para 747), com finais menores, e só
voltou ao nível de 2023 em 2026 (863). O enxame (Figura~\ref{fig:enxame}) mostra cada resultado por prova e
competição: a mediana do IT é 529 no InterUSP e 547 no TUSCA, uma diferença pequena diante da dispersão dentro de
cada competição.

\begin{figure}[H]
\centering
\anotada[.8\linewidth]{obs-ridge.png}{\destaque{0.005}{0.07}{0.40}{0.135}{1}\destaque{0.535}{0.255}{0.595}{0.61}{2}\destaque{0.005}{0.70}{0.995}{0.79}{3}}
\caption{Nível (Q1): densidade do Índice Técnico por edição, InterUSP em azul e TUSCA em amarelo. \marca{1} Filtros de competição, ano, sexo e prova. \marca{2} A linha tracejada liga as medianas do InterUSP: 556, 512 e 524. \marca{3} TUSCA 2024: a distribuição inteira desce, inclusive o topo.}
\label{fig:nivel}
\end{figure}

\begin{figure}[H]
\centering
\anotada[.8\linewidth]{obs-swarm.png}{\destaque{0.005}{0.225}{0.995}{0.285}{1}\destaque{0.31}{0.075}{0.50}{0.105}{2}\destaque{0.86}{0.30}{0.995}{0.35}{3}}
\caption{Enxame (Q2): cada ponto é um resultado, colorido pela competição. \marca{1} Uma linha por prova, com os resultados pelo IT. \marca{2} Modo de cor (competição, equipe ou estilo). \marca{3} Os melhores resultados, com IT perto de 1.000. Passar o mouse mostra atleta e tempo; clicar abre a ficha.}
\label{fig:enxame}
\end{figure}

\paragraph{Q3: equipes.} O \emph{bump chart} (Figura~\ref{fig:ranking}) mostra que o InterUSP alterna entre POLI
(campeã em 2023 e 2026) e MED PINHEIROS (2024 e 2025), sempre com MED RIBEIRÃO em terceiro, enquanto o TUSCA é
mais disputado: POLI em 2023, UFSCAR em 2024 e 2025 e LAP em 2026, saindo do meio da tabela. A matriz equipe
$\times$ prova (Figura~\ref{fig:matriz}) mostra a fatia dos pontos de cada prova que cada equipe conquistou;
ordenada por similaridade, ela revela blocos de especialização: a CAASO soma 35\% dos pontos do 200 m livre e
38\% do 400 m livre, enquanto a POLI lidera três das quatro provas de 50 m (no 50 m livre fica a dois pontos da CAASO), o 100 m medley e o 4$\times$50 m livre, e a
UFSCAR aparece entre as três primeiras em quase todas as provas sem liderar nenhuma isoladamente.

\begin{figure}[H]
\centering
\begin{subfigure}{.49\linewidth}\anotada{obs-bump.png}{\destaque{0.07}{0.21}{0.95}{0.30}{1}\destaque{0.07}{0.315}{0.95}{0.345}{2}\destaque{0.07}{0.605}{0.15}{0.64}{3}\destaque{0.60}{0.60}{0.93}{0.80}{4}}\caption{Ranking por edição. \marca{1} InterUSP: POLI e MED PINHEIROS alternam o 1º lugar. \marca{2} MED RIBEIRÃO sempre em 3º. \marca{3} TUSCA 2023: POLI campeã. \marca{4} LAP sobe do 5º ao 1º lugar em 2026.}\label{fig:ranking}\end{subfigure}\hfill
\begin{subfigure}{.49\linewidth}\anotada{obs-matrix.png}{\destaque{0.185}{0.16}{0.29}{0.19}{1}\destaque{0.62}{0.27}{0.80}{0.295}{2}\destaque{0.18}{0.30}{0.45}{0.33}{3}}\caption{Pontos por equipe e prova. \marca{1} Ordem pelo total ou por similaridade. \marca{2} CAASO lidera o 200 e o 400 m livre. \marca{3} POLI lidera o 50 m costas, o 50 m peito e o 50 m borboleta.}\label{fig:matriz}\end{subfigure}
\caption{Equipes (Q3).}
\end{figure}

\paragraph{Q4: fases da prova.} As súmulas não têm parciais, então o Raio-X estima as fases (reação, voo,
submerso, nado, viradas e chegada) com um modelo calibrado por médias de finalistas de elite~\cite{arellano}, que escala cada
fase pelo nível do nadador (tempo de referência $\div$ tempo); esses valores aparecem marcados como estimados na
interface. A fita (Figura~\ref{fig:raiox}a) mostra as fases de cada finalista lado a lado, e o modo ``todas as
provas'' (Figura~\ref{fig:raiox}b) compara cada métrica entre todas as finais. Somando todas as provas
individuais, a perda de tempo para a referência de elite fica em 57\% no nado, 16\% na saída e 21\% nas viradas.
A Arena 3D (Figura~\ref{fig:arena}) anima a final escolhida com o visitante numa das raias, com câmeras de
transmissão, placar e telemetria ao vivo.

\begin{figure}[H]
\centering
\begin{subfigure}{.49\linewidth}\anotada{raiox-fita.png}{\destaque{0.005}{0.165}{0.995}{0.225}{1}\destaque{0.10}{0.325}{0.83}{0.355}{2}\destaque{0.10}{0.555}{0.94}{0.585}{3}}\caption{Fita de uma final. \marca{1} Telemetria do nadador escolhido. \marca{2} Referência de elite (PRO). \marca{3} As fases de cada finalista: reação, voo, submerso e nado.}\end{subfigure}\hfill
\begin{subfigure}{.49\linewidth}\anotada{raiox-todas.png}{\destaque{0.635}{0.12}{0.69}{0.15}{1}\destaque{0.015}{0.365}{0.045}{0.405}{2}\destaque{0.052}{0.21}{0.072}{0.55}{3}}\caption{Saída até 15 m em todas as provas. \marca{1} Métrica comparada. \marca{2} Referência PRO. \marca{3} Mediana da prova.}\end{subfigure}
\caption{Raio-X (Q4).}
\label{fig:raiox}
\end{figure}

\begin{figure}[H]
\centering
\begin{subfigure}{.49\linewidth}\anotada{arena.png}{\destaque{0.01}{0.085}{0.29}{0.125}{1}\destaque{0.005}{0.185}{0.18}{0.47}{2}\destaque{0.80}{0.085}{0.995}{0.41}{3}\destaque{0.245}{0.935}{0.755}{0.985}{4}}\caption{Arena 3D. \marca{1} Placar e cronômetro. \marca{2} Telemetria ao vivo. \marca{3} Classificação parcial. \marca{4} Câmeras e controles.}\label{fig:arena}\end{subfigure}\hfill
\begin{subfigure}{.49\linewidth}\anotada{ficha-atleta.png}{\destaque{0.725}{0.08}{0.995}{0.27}{1}\destaque{0.725}{0.28}{0.995}{0.52}{2}\destaque{0.725}{0.54}{0.995}{0.79}{3}}\caption{Ficha do atleta. \marca{1} Medalhas e pontos. \marca{2} Trajetória do IT por ano e competição. \marca{3} Perfil por estilo.}\label{fig:ficha}\end{subfigure}
\caption{A experiência e o detalhe sob demanda.}
\end{figure}

\section{Discussões preliminares}\label{sec:disc}

\begin{itemize}[leftmargin=1.2em]
  \item \textbf{Q1.} Não há tendência de alta na mediana: ela cai até 2025 e se recupera em 2026 nas duas
  competições. No InterUSP, porém, os melhores ficaram mais rápidos enquanto as finais cresceram; no TUSCA, 2024
  foi uma edição mais fraca em toda a distribuição.
  \item \textbf{Q2.} O TUSCA tem mediana um pouco maior (547 contra 529), mas as distribuições se sobrepõem
  quase inteiras; a diferença é menor que a variação entre edições da mesma competição.
  \item \textbf{Q3.} O InterUSP é um duelo POLI $\times$ MED PINHEIROS; o TUSCA teve três campeões em quatro anos.
  As equipes se especializam: CAASO no fundo (200 e 400 m), POLI na velocidade (50 m).
  \item \textbf{Q4.} A maior parte da diferença para a elite está no nado, mas saída e viradas somam mais de um
  terço, uma parte treinável que a súmula sozinha não mostra.
\end{itemize}

\paragraph{Limitações.} As fases da prova e as braçadas são estimativas de um modelo, não medições. O IT usa a
melhor marca da base como referência, então mede o nível relativo a estas competições, não o absoluto. A ligação de
atletas é heurística e pode juntar homônimos ou separar a mesma pessoa com grafias muito diferentes. As súmulas
não informam o tamanho da piscina.

\paragraph{Lições.} A maior parte do esforço foi o pré-processamento: cada edição tem um layout de súmula diferente,
e erros pequenos (uma colocação digitada errada, um CPF colado ao nome da equipe) só aparecem quando os totais
são conferidos com a classificação geral. Separar a experiência (piscina 3D) da análise (visões em d3) ajudou a
manter cada visão simples e ligada a uma questão.

\section{Próximos passos}\label{sec:prox}

A segunda etapa vai aprofundar a análise de atletas e equipes com duas novas questões: existem perfis de atletas
(velocistas, fundistas, especialistas em um estilo)? Os pontos de cada equipe dependem de poucos atletas? Para
isso, entram novas visualizações coordenadas às atuais e novos mecanismos de interação entre as visões, além de
uma estimativa das chances de um nadador numa final a partir do seu tempo.

\begin{thebibliography}{99}\small
\bibitem{eban-interusp} EBAN Swim. Resultados do InterUSP 2023, 2024, 2025 e 2026.
\url{https://ebanswim.com.br/interusp-2023/} (e páginas de 2024 a 2026). Acesso em nov.~2026.
\bibitem{eban-tusca} EBAN Swim. Resultados do TUSCA 2023, 2024, 2025 e 2026.
\url{https://ebanswim.com.br/tusca-2023/}, \url{https://ebanswim.com.br/tusca-2024/},
\url{https://ebanswim.com.br/tusca-2025-2/}, \url{https://ebanswim.com.br/tusca-central-2026/}. Acesso em nov.~2026.
\bibitem{wa-regras} World Aquatics. \emph{Competition Regulations: Swimming Rules (SW)}. \url{https://www.worldaquatics.com/swimming/rules}. Acesso em nov.~2026.
\bibitem{cbda} Confederação Brasileira de Desportos Aquáticos (CBDA). \emph{Regras oficiais de natação} e regulamentos de competição. \url{https://www.cbda.org.br}. Acesso em nov.~2026.
\bibitem{fap} Federação Aquática Paulista (FAP). \url{https://www.aquaticapaulista.org.br}. Acesso em nov.~2026.
\bibitem{cbdu} Confederação Brasileira do Desporto Universitário (CBDU). Jogos Universitários Brasileiros (JUBs). \url{https://www.cbdu.org.br}. Acesso em nov.~2026.
\bibitem{arellano} R. Arellano, P. Brown, J. Cappaert, R. C. Nelson. Analysis of 50-, 100-, and 200-m freestyle swimmers at the 1992 Olympic Games. \emph{Journal of Applied Biomechanics}, 10(2):189--199, 1994.
\bibitem{wa-points} World Aquatics. \emph{World Aquatics Points}. \url{https://www.worldaquatics.com/swimming/points}.
\bibitem{d3} M. Bostock, V. Ogievetsky, J. Heer. D$^3$: Data-Driven Documents. \emph{IEEE TVCG}, 17(12):2301--2309, 2011.
\bibitem{threejs} three.js --- JavaScript 3D library. \url{https://threejs.org}.
\bibitem{shneiderman} B. Shneiderman. The eyes have it: a task by data type taxonomy for information
visualizations. \emph{IEEE Symposium on Visual Languages}, p.~336--343, 1996.
\bibitem{okabe-ito} M. Okabe, K. Ito. \emph{Color Universal Design (CUD): how to make figures and presentations that
are friendly to colorblind people}. 2008. \url{https://jfly.uni-koeln.de/color/}.
\bibitem{behrisch} M. Behrisch et al. Matrix reordering methods for table and network visualization.
\emph{Computer Graphics Forum}, 35(3):693--716, 2016.
\bibitem{munzner} T. Munzner. \emph{Visualization Analysis and Design}. CRC Press, 2014.
\end{thebibliography}

\end{document}
